\section{Rocket 2D}\label{sec:rocket2d}
\autoref{sec:notation} lists every symbol used in this report together with the identifier the
simulator uses for it.

\subsection{Frames}\label{sec:frames}
\begin{itemize}
    \item $\mathcal{F}_0$ --- ground/launch frame, origin at the launch point $O^*$.
    \item $\mathcal{F}_g$ --- ground/gravity frame: $\hat{x}_g$ downrange, $\hat{z}_g$ down, along
    gravity.
    \item $\mathcal{F}_r$ --- rocket (body) frame: $\hat{x}_r$ along the nose axis, forward.
    \item $\mathcal{F}_t$ --- thrust frame: $\hat{x}_t$ along the line of thrust.
\end{itemize}
$\mathcal{F}_0$ is fixed at the launch point $O^*$, with $\hat{x}_0$ horizontal (downrange) and
$\hat{z}_0$ down. The rocket's state is always integrated in this frame, regardless of which gravity
or ground model is switched on --- it is the one fixed thing in the whole document. The
translational state is the position $\vec{x}_r=x\,\hat{x}_0+z\,\hat{z}_0$, the velocity
$\vec{v}_r=\mathrm{d}\vec{x}_r/\mathrm{d}t$ and the acceleration
$\vec{a}_r=\mathrm{d}\vec{v}_r/\mathrm{d}t$, all resolved in $\mathcal{F}_0$.

$\mathcal{F}_g$ carries the direction gravity actually points, $\hat{z}_g(\varphi)$, a function of
the downrange arc angle $\varphi$ introduced in \autoref{sec:gravity}. It coincides with
$\mathcal{F}_0$ by default, $\hat{x}_g(0)\equiv\hat{x}_0$ and $\hat{z}_g(0)\equiv\hat{z}_0$, and
only rotates away from it once a non-parallel direction model is enabled --- so $\mathcal{F}_g$
and $\mathcal{F}_0$ are kept notationally separate precisely because ``ground'' and ``the
direction gravity points'' are not always the same thing, even though they agree at the launch
point.

$\mathcal{F}_r$ is right-handed, with $\hat{x}_r$ along the rocket's nose axis (positive forward),
$\hat{z}_r$ completing the pitch plane with it, and $\hat{y}_r$ out of the pitch plane. The pitch
angle $\theta_r$ is measured from $\hat{x}_0$ (the ground-frame horizontal), positive nose-up; the
pitch rate is $\omega_r=\mathrm{d}\theta_r/\mathrm{d}t$, the rotation rate about $\hat{y}_r$, and
the angular acceleration is $\varepsilon_r=\mathrm{d}\omega_r/\mathrm{d}t$. In ground-frame
components,
\begin{equation}\label{eq:body_axes_in_ground}
    \hat{x}_r = \cos\theta_r\,\hat{x}_0 - \sin\theta_r\,\hat{z}_0, \qquad
    \hat{z}_r = \sin\theta_r\,\hat{x}_0 + \cos\theta_r\,\hat{z}_0.
\end{equation}
($\alpha$ is reserved for the angle of attack, \autoref{sec:relwind}.)

$\mathcal{F}_t$ carries the actual line of thrust action, $\hat{x}_t$, which need not coincide
with the nose axis: a real nozzle can be canted, or otherwise misaligned, by a small angle
$\theta_t$ from the body axis --- a fixed property of how the motor is mounted, not of the flight
state. Giving thrust its own frame, rather than folding $\theta_t$ into every thrust term by hand,
is what lets that distinction be written down once:
\begin{equation}\label{eq:thrust_frame}
    \hat{x}_t = \cos\theta_t\,\hat{x}_r - \sin\theta_t\,\hat{z}_r, \qquad
    \hat{z}_t = \sin\theta_t\,\hat{x}_r + \cos\theta_t\,\hat{z}_r, \qquad \hat{y}_t\equiv\hat{y}_r.
\end{equation}
$\mathcal{F}_t$ is obtained from $\mathcal{F}_r$ exactly as $\mathcal{F}_r$ is from
$\mathcal{F}_0$ in \autoref{eq:body_axes_in_ground}, so a positive $\theta_t$ tilts the thrust line
in the same sense as a positive $\theta_r$ tilts the nose. The older cant angle $\delta$ had the
opposite sign, $\theta_t=-\delta$ (\autoref{sec:notation}). The thrust direction can in principle
be controlled; for now $\theta_t$ is carried as a constant input, ideally zero --- a fixed
manufacturing-type misalignment, swept by hand to see how sensitive the airframe is to it. A
commanded, time-varying gimbal angle driven by a control law is future work
(\autoref{sec:backlog}).

\subsection{Forces and moments}\label{sec:forces}
Weight is the main force the rocket has to overcome to increase its altitude:
\begin{equation}
    \vec{W} = mg\hat{z}_g,
\end{equation}
where $m$ is the rocket's mass and $g$ the gravitational acceleration along the local vertical
$\hat{z}_g$. Both factors are models rather than constants, and are chosen independently: the
magnitude $g$ and the direction $\hat{z}_g$ answer different questions --- how strong, and which
way --- and are supplied separately by \autoref{sec:gravity}.

Thrust is generated by ejecting propellant at high velocity, plus a contribution from the pressure
difference across the nozzle exit. With the thrust frame of \autoref{eq:thrust_frame} it is simply
\begin{equation}
    \vec{T} = T\hat{x}_t = T\left(\cos\theta_t\,\hat{x}_r - \sin\theta_t\,\hat{z}_r\right),
\end{equation}
purely axial in its own frame, with $T$ the thrust magnitude --- the measured thrust curve plus a
nozzle-exit pressure correction (\autoref{sec:thrust}); the angle $\theta_t$ inside $\hat{x}_t$
carries the whole of the misalignment, so the first form never needs to be rewritten as $\theta_t$
changes.

Motion through the surrounding air produces an aerodynamic force, and, about a reference point $p$
on the rocket (\autoref{sec:geometry}), an aerodynamic pitching moment:
\begin{equation}
    \vec{A}_f = -qS\left(C_x\hat{x}_r + C_z\hat{z}_r\right), \qquad
    \vec{A}_m = -qSc\,C_m\hat{y}_r, \qquad
    q = \tfrac12\rho v_{rel}^2,
\end{equation}
with $q$ the dynamic pressure, $\rho$ the air density, $v_{rel}$ the airspeed (the magnitude of
the rocket's velocity relative to the air, \autoref{sec:relwind}), $S$ the aerodynamic reference
area and $c$ the reference length. $C_x,C_z,C_m$ are the axial-force, normal-force and
pitching-moment coefficients, defined in the body frame (\autoref{sec:aero}). The aerodynamic
force and moment are \emph{always} these dynamic-pressure-scaled coefficient products, whichever
of \autoref{sec:aero}'s coefficient models is selected; no model in this document supplies an
aerodynamic force directly as a force.

\subsection{Geometry}\label{sec:geometry}
Each of the three forces acts at its own point on the rocket:
\begin{itemize}
    \item $A$ --- nose tip, the origin of every length below;
    \item $g$ --- centre of gravity, where weight acts;
    \item $t$ --- thrust application point, where the (possibly canted) thrust line of
    $\mathcal{F}_t$ crosses the body;
    \item $p$ --- aerodynamic reference point (centre of pressure), where $\vec{A}_f$ acts and
    about which $\vec{A}_m$ is taken.
\end{itemize}
Every dimension is defined in the rocket frame $\mathcal{F}_r$: $l$ is the axial distance of a
point from $A$, measured aft (positive towards the tail), and $r$ its radial distance from the
axis, along $-\hat{z}_r$. Writing $(a_1,a_2,a_3)_r\equiv a_1\hat{x}_r+a_2\hat{y}_r+a_3\hat{z}_r$
for components in $\mathcal{F}_r$,
\begin{equation}
    \vec{Ag} = -l_g\hat{x}_r - r_g\hat{z}_r, \qquad
    \vec{At} = -l_t\hat{x}_r - r_t\hat{z}_r, \qquad
    \vec{Ap} = -l_p\hat{x}_r.
\end{equation}
The moments below are taken about $g$, so what they need is the position of the other two points
relative to it:
\begin{equation}\label{eq:arms}
    \vec{gt} = -\vec{Ag} + \vec{At} = (l_g-l_t,\;0,\;r_g-r_t)_r, \qquad
    \vec{gp} = -\vec{Ag} + \vec{Ap} = (l_g-l_p,\;0,\;r_g)_r.
\end{equation}

The radial distances $r$ describe the axis misalignment: in a real rocket the centre of gravity
and the nozzle cannot be strictly on the axis, whether from manufacturing tolerance or from an
as-built mass asymmetry. $r$ can be zero, and is by default. As $p$ is a reference point rather
than a measured one, no misalignment is defined for it ($r_p\equiv0$); any asymmetry in the
external surfaces instead shows up in the aerodynamic coefficients themselves, e.g.\ in polar
plots, not in a radial offset of $p$. $r_g$ and $r_t$ are carried as ordinary constants alongside
$l_g,l_t$ throughout what follows, on the same footing as the thrust angle $\theta_t$ --- setting
them to zero recovers the on-axis geometry exactly, and there is no separate ``model'' for them
beyond that: like $\theta_t$, they are inputs representing an as-built offset, in principle
obtained from manufacturing metrology of the actual airframe rather than derived from anything
else in this document. $l_t$, by contrast, only matters when the thrust is canted: it enters the
moment through $\sin\theta_t$ alone (\autoref{eq:eom_rot_expanded}), so for $\theta_t=0$ its value
is immaterial and it can be set to zero.

\subsection{Equations of motion}\label{sec:eom}
\begin{gather}
    m\vec{a}_r = \sum \vec{F} = \vec{W} + \vec{T} + \vec{A}_f \label{eq:eom_trans}\\
    I\varepsilon_r = \sum \vec{M}(g) \cdot \hat{y}_r
    = \left[\vec{gt}\times\vec{T} + \vec{gp}\times\vec{A}_f + \vec{A}_m\right]\cdot \hat{y}_r
    \label{eq:eom_rot}
\end{gather}
Substituting the forces of \autoref{sec:forces},
\begin{equation}\label{eq:eom_trans_expanded}
    m\vec{a}_r = mg\hat{z}_g + T\left(\cos\theta_t\,\hat{x}_r - \sin\theta_t\,\hat{z}_r\right)
    - qS\left(C_x\hat{x}_r + C_z\hat{z}_r\right).
\end{equation}
For the rotation, substitute \autoref{eq:arms} and use the $\hat{y}_r$ component of the cross
product of two pitch-plane vectors,
$\left[(a_1,0,a_3)_r\times(b_1,0,b_3)_r\right]\cdot\hat{y}_r = a_3b_1-a_1b_3$:
\begin{align}
    I\varepsilon_r
    &= \left[T\left(-\vec{Ag}+\vec{At}\right)\times\hat{x}_t
       + \left(-\vec{Ag}+\vec{Ap}\right)\times\vec{A}_f + \vec{A}_m\right]\cdot\hat{y}_r
       \nonumber\\
    &= \Bigl[T\,(l_g-l_t,\,0,\,r_g-r_t)_r\times(\cos\theta_t,\,0,\,-\sin\theta_t)_r
       \nonumber\\
    &\qquad + qS\,(l_g-l_p,\,0,\,r_g)_r\times(-C_x,\,0,\,-C_z)_r
       - qSc\,C_m\hat{y}_r\Bigr]\cdot\hat{y}_r
       \nonumber\\
    &= T\left[(r_g-r_t)\cos\theta_t + (l_g-l_t)\sin\theta_t\right]
       + qS\left[-r_gC_x + (l_g-l_p)C_z - c\,C_m\right]. \label{eq:eom_rot_expanded}
\end{align}
Each term reads off the geometry directly. The radial offset $r_g-r_t$ of the thrust point gives
an arm to the axial thrust component $T\cos\theta_t$, and its axial offset $l_g-l_t$ an arm to the
normal component $T\sin\theta_t$; in the same way $r_g$ is the arm of the axial aerodynamic force
$qSC_x$ and $l_g-l_p$ that of the normal one $qSC_z$. For the aligned rocket,
$r_g=r_t=\theta_t=0$, thrust drops out of the moment altogether and
$I\varepsilon_r=qS\left[(l_g-l_p)C_z-c\,C_m\right]$.

Moments are taken about the centre of gravity $g$, so weight contributes none:
$\vec{gg}\times\vec{W}=0$, for whatever direction $\vec{W}$ points. $m$ and $I$ are not constants
once the propellant model of \autoref{sec:mass} is active --- they become $m(t), I(t)$, and the
centre of gravity $g$ itself travels as $l_g(t)$. Reinterpreting $\mathcal{F}_0$ as Earth-fixed
rather than inertial adds one further translational term to \autoref{eq:eom_trans} and nothing at
all to \autoref{eq:eom_rot} (\autoref{sec:noninertial}). \autoref{sec:analysis} expands these two
equations by hand.

These two equations describe \emph{free} flight, which is all of a flight except its two
mechanically constrained phases. While the rocket is still on the launch rail
(\autoref{sec:rail}) and while it hangs under a deployed recovery canopy (\autoref{sec:recovery}),
the attitude state is frozen, $\omega_r=\varepsilon_r=0$, and \autoref{eq:eom_trans} is modified
--- projected onto the rail direction in the first case, given one additional inertial-frame force
in the second. Both are written up with the models rather than here, because each needs its own
constraint or force model to be stated first; the equations above are what they modify, and what
they revert to.

\subsubsection{Model inventory}\label{sec:eom_inventory}
Every symbol on the right-hand sides of \autoref{eq:eom_trans_expanded} and
\autoref{eq:eom_rot_expanded} is either a state or is supplied by a model:
\begin{itemize}
    \item $m$, $I$ --- mass and moment of inertia, from the mass model (\autoref{sec:mass});
    \item $g\hat{z}_g$ --- gravity, magnitude and direction (\autoref{sec:gravity});
    \item $T$ --- thrust curve plus pressure correction (\autoref{sec:thrust});
    \item $\theta_t$ --- thrust direction: can be controlled, for now a constant, ideally zero
    (\autoref{sec:thrust});
    \item $q$ --- from the atmosphere model, the state and the wind model, through $\rho$ and
    $v_{rel}$ (\autoref{sec:atmosphere}, \autoref{sec:wind}, \autoref{sec:relwind});
    \item $S$, $c$ --- reference area and reference length (\autoref{sec:aero});
    \item $C_x$, $C_z$, $C_m$ --- aerodynamic coefficients (\autoref{sec:aero});
    \item $l_t$, $r_t$, $l_p$, $l_g$, $r_g$ --- lengths (\autoref{sec:geometry}); $l_g(t)$ moves
    with the mass model (\autoref{sec:mass}), and the position of $p$ is supplied together with the
    coefficients (\autoref{sec:aero});
    \item states --- $\theta_r$, $\omega_r$, $\varepsilon_r$ and $\vec{x}_r$, $\vec{v}_r$,
    $\vec{a}_r$ (\autoref{sec:frames}); the equations deliver $\varepsilon_r$ and $\vec{a}_r$, the
    rest is integrated;
    \item and, outside free flight, the parachute (\autoref{sec:recovery}) and the launch rail
    (\autoref{sec:rail}).
\end{itemize}

\subsection{Notation}\label{sec:notation}
The table maps each symbol of this report to the identifier the simulator uses for it
(\texttt{sim/} keeps its own names) and to the symbol used in earlier versions of this report.
``Not stored'' means the code keeps only an arm relative to the centre of gravity, not the length
from the nose tip; ``not implemented'' means the quantity appears in the report but not in the
simulator.
{\small
\begin{longtable}{>{\raggedright\arraybackslash}p{2.6cm}
                >{\raggedright\arraybackslash}p{2.0cm}
                >{\raggedright\arraybackslash}p{4.1cm}
                >{\raggedright\arraybackslash}p{4.4cm}
                >{\raggedright\arraybackslash}p{1.5cm}}
\hline
\textbf{Report} & \textbf{Earlier doc} & \textbf{Code} & \textbf{Meaning} & \textbf{Units} \\
\hline
\endfirsthead
\hline
\textbf{Report} & \textbf{Earlier doc} & \textbf{Code} & \textbf{Meaning} & \textbf{Units} \\
\hline
\endhead
\hline
\endfoot
\multicolumn{5}{l}{\textbf{Geometry and frames}} \\
\hline
$A$ & $A$ & --- & nose tip, origin of all lengths & --- \\
$g$ & $B$ & --- & centre of gravity & --- \\
$t$ & $C$ & --- & thrust application point & --- \\
$p$ & $D$ & --- & aerodynamic reference point (centre of pressure) & --- \\
$l_g$ & $l_B$ & \texttt{mass.cog\_position(m\_p, p)}; local \texttt{l\_B} in \texttt{sim/dynamics.py} (mass model only) & axial distance $A\to g$ & m \\
$r_g$ & $u_B$ & \texttt{p.lengths.u\_B} & radial offset of $g$ & m \\
$l_t$ & $l_C$ & not stored; $l_t-l_g$ = \texttt{p.lengths.l1 + p.lengths.l2} & axial distance $A\to t$ & m \\
$r_t$ & $u_C$ & \texttt{p.lengths.u\_C} & radial offset of $t$ & m \\
$l_p$ & $l_D$ & not stored; $l_p-l_g$ = \texttt{l1} (row below) & axial distance $A\to p$ & m \\
$r_p\equiv0$ & $u_D\equiv0$ & --- & no radial offset of $p$, by definition & m \\
$l_p-l_g$ & $l_1=l_D-l_B$ & \texttt{p.lengths.l1}; \texttt{aero.coefficients(\dots).l1} (varies under the table model) & arm $g\to p$, positive with $p$ aft & m \\
$l_t-l_p$ & $l_2=l_C-l_D$ & \texttt{p.lengths.l2} & arm $p\to t$ & m \\
$\vec{gt}$, $\vec{gp}$ & $\vec{BC}$, $\vec{BD}$ & \texttt{aB}, \texttt{aA} in \texttt{sim/dynamics.py} & position of $t$, $p$ relative to $g$ & m \\
$\theta_t$ & $-\delta$ & $-$\texttt{np.radians(}\allowbreak\texttt{p.lengths.misalign\_deg)} & thrust-line angle from $\hat{x}_r$ & rad (code: deg) \\
\hline
\multicolumn{5}{l}{\textbf{Mass}} \\
\hline
$m$, $I$ & $m$, $I$ & \texttt{p.weight.m}, \texttt{p.weight.I}, or \texttt{mass.moment\_of\_inertia} & mass; pitch moment of inertia about $g$ & kg, $\mathrm{kg\,m^2}$ \\
$m_r$ & $m_d$ less the casing & \texttt{airframe.m} (\texttt{configs/schema.py}); \texttt{p.weight.m\_d} also includes the casing & dry-rocket mass (constant) & kg \\
$m_m(t)$ & $m_p$ (propellant only) & \texttt{mass.propellant\_mass(m, p)} (\texttt{y[M]} $-$ \texttt{m\_d}) is the propellant only; $m_m$ adds the casing $m_f$ & motor mass (casing and propellant), burned down to $m_f$ & kg \\
$m_0$, $m_f$ & --- & \texttt{motor.m\_total}; $m_f$ = \texttt{m\_total - m\_prop} & motor mass at ignition, at burnout & kg \\
$\Delta t$ & --- & not implemented (no linear mass model) & burn duration (linear mass model) & s \\
$\Delta m$ & --- & \texttt{motor.m\_prop} & burned mass $m_0-m_f$ & kg \\
$\mathrm{ISP}$ & $I_{sp}$ & \texttt{p.weight.I\_sp}; \texttt{motor.I\_sp} & specific impulse & s \\
$g_0$ & $g_0$ & \texttt{p.g0} & standard gravity & $\mathrm{m/s^2}$ \\
$I_r$, $I_{r,g}$ & $I_d$ & \texttt{p.weight.I\_d} (dry assembly about its own centroid) & dry-rocket inertia about $g$; about its own centre of gravity & $\mathrm{kg\,m^2}$ \\
$I_m(t)$, $I_{m,g}(t)$ & $I_p(t)$ & \texttt{p.weight.I\_p\_spec}$\cdot m_p$ (propellant only) & motor inertia about $g$; about its own centre of gravity & $\mathrm{kg\,m^2}$ \\
$I_0$, $I_f$ & --- & \texttt{motor.I\_0}, \texttt{motor.I\_f} & $I_{m,g}$ at ignition, at burnout & $\mathrm{kg\,m^2}$ \\
$l_{g,r}$, $r_{g,r}$ & $l_d$ & \texttt{p.weight.l\_d}; radial: not stored (constant \texttt{u\_B}) & dry-rocket centre of gravity (axial, radial) & m \\
$l_{g,m}(t)$, $r_{g,m}(t)$ & $l_p$ (propellant) & \texttt{p.weight.l\_p} (fixed); radial: not stored & motor centre of gravity (axial, radial) & m \\
$l_0$, $l_f$ & --- & \texttt{motor.l\_cog\_0}, \texttt{motor.l\_cog\_f} & $l_{g,m}$ at ignition, at burnout & m \\
\hline
\multicolumn{5}{l}{\textbf{Forces and aerodynamics}} \\
\hline
$S$ & $S$ & \texttt{p.aero.S} & aerodynamic reference area (same symbol in both) & $\mathrm{m^2}$ \\
$c$ & $d_{ref}$ & \texttt{p.aero.d\_ref} & reference length of the pitching moment & m \\
$C_x$, $C_z$, $C_m$ & $C_x$, $C_z$, $C_m$ ($C_A$, $C_N$) & \texttt{p.aero.C\_x}, \texttt{C\_z}, \texttt{C\_m} (constant model); fields of \texttt{aero.AeroCoeffs} & axial-force, normal-force, pitching-moment coefficients & --- \\
$qSC_x$, $qSC_z$ & $A_x$, $A_z$ & \texttt{forces.aero\_force\_body} returns $(-A_x,-A_z)$ & aerodynamic force components & N \\
$qSc\,C_m$ & $A_m$ & \texttt{forces.aero\_moment\_y} returns $-A_m$ & aerodynamic pitching moment & N\,m \\
$S_k$, $c_k$, $C_{x,k}$, $C_{z,k}$, $C_{m,k}$, $pp_k$ & --- & not implemented & aero component $k$: reference area and length, coefficients, reference point & $\mathrm{m^2}$, m, ---, m \\
$C_{x,0}$, $k$ & $c_{x,0}$, $1/(e\pi\lambda)$ & not implemented & drag polar $C_x=C_{x,0}+kC_z^2$ & --- \\
$\partial C/\partial\alpha$, $\partial C/\partial M$, $\partial C/\partial Re$ & --- & only the $\alpha$ slope is in the code & linearized-model slopes & 1/rad, ---, --- \\
$C_{z,\alpha}$, $C_{m,\alpha}$, $C_{m,q}$, $C_{z,0}$, $C_{m,0}$ & $c_{z,\alpha}$, $c_{m,\alpha}$, $c_{m,q}$, \dots & aero coefficient inputs (\texttt{c\_mq} column of the aero table, \dots) & coefficient slopes and offsets & 1/rad or --- \\
\hline
\multicolumn{5}{l}{\textbf{Thrust}} \\
\hline
$T$, $\vec{T}$ & $T$, $\vec{T}$ & \texttt{thrust\_curve.thrust}; \texttt{forces.thrust\_body} & thrust magnitude; vector in $\mathcal{F}_r$ & N \\
$T_0(t)$ & $T_{ref}(t)$ & \texttt{thrust\_curve} (\texttt{T\_ref}) & reference (test-stand) thrust & N \\
$p_0$ & $p_{ref}$ & \texttt{p.thrust.p\_ref} & ambient pressure during the static test & Pa \\
$A_e$ & $A_e$ & \texttt{p.thrust.A\_e} & nozzle exit area & $\mathrm{m^2}$ \\
\hline
\multicolumn{5}{l}{\textbf{Atmosphere and wind}} \\
\hline
$p_{atm}$ & $p_{amb}$ & \texttt{air\_state(\dots).p} & ambient pressure & Pa \\
$q$ & $q$ & \texttt{flow.q} & dynamic pressure & Pa \\
$v_{rel}$, $\vec{v}_{rel}$ & $V$ & \texttt{flow.V} & airspeed; relative-wind vector & m/s \\
$v_{rel,x}$, $v_{rel,z}$ & $u$, $w$ & --- & relative-velocity components in $\mathcal{F}_r$ & m/s \\
$\alpha$ & $\alpha$ & \texttt{flow.alpha} & angle of attack & rad \\
$w_{up}$ & --- & not implemented & updraft speed ($w_z=-w_{up}$) & m/s \\
$v_{ref}$ & $V_{ref}$ & turbulence setting & fixed reference airspeed of the Dryden spectrum & m/s \\
\hline
\multicolumn{5}{l}{\textbf{States and integration}} \\
\hline
$\vec{y}=(x,z,v_x,$ $v_z,\theta_r,\omega_r,m)$ & $\vec{y}$ & \texttt{y[X,}\allowbreak\texttt{Z,}\allowbreak\texttt{VX,}\allowbreak\texttt{VZ,}\allowbreak\texttt{THETA,}\allowbreak\texttt{OMEGA,}\allowbreak\texttt{M]} & integrated state vector & m, m/s, rad, rad/s, kg \\
$\theta_r$ & $\theta$ & \texttt{y[THETA]} & pitch angle & rad \\
$\omega_r$ & $\omega$ & \texttt{y[OMEGA]} & pitch rate & rad/s \\
$\varepsilon_r$ & $\mathrm{d}\omega/\mathrm{d}t$ & \texttt{M\_y / I} ($\mathrm{d}\omega/\mathrm{d}t$ in \texttt{sim/dynamics.py}) & angular acceleration & $\mathrm{rad/s^2}$ \\
$\vec{x}_r$ & $\vec{x}=(x,z)$ & \texttt{y[X]}, \texttt{y[Z]} & position in $\mathcal{F}_0$ & m \\
$\vec{v}_r$ & $\mathrm{d}\vec{x}/\mathrm{d}t$ & \texttt{y[VX]}, \texttt{y[VZ]} & velocity in $\mathcal{F}_0$ & m/s \\
$\vec{v}_g$ & $\vec{v}_B$ & \texttt{y[VX]}, \texttt{y[VZ]} & CoG velocity (the same vector as $\vec{v}_r$) & m/s \\
$v_x$, $v_z$ & $\mathrm{d}x/\mathrm{d}t$, $\mathrm{d}z/\mathrm{d}t$ & \texttt{y[VX]}, \texttt{y[VZ]} & velocity components in $\mathcal{F}_0$ & m/s \\
$\vec{a}_r$ & $\mathrm{d}\vec{v}/\mathrm{d}t$ & \texttt{a\_i} in \texttt{sim/dynamics.py} & acceleration in $\mathcal{F}_0$ & $\mathrm{m/s^2}$ \\
$m_{launch}$ & --- & \texttt{s0.m}; \texttt{p.weight.m} & liftoff mass $m(0)$ & kg \\
$h$ & $dt$ & \texttt{p.solver.dt} & fixed RK4 step & s \\
$\theta_{rail}$, $L_{rail}$, $s_{rail}$, $\hat{e}_{rail}$ & same & \texttt{p.site.rail\_angle\_deg}, \texttt{p.site.L\_rail}, \texttt{rail\_distance}, \texttt{\_e\_rail} & rail elevation, length, distance travelled, direction & rad, m, m, --- \\
\hline
\multicolumn{5}{l}{\textbf{Recovery and air brake}} \\
\hline
$S_{ch}$, $C_{D,ch}$, $\vec{F}_{ch}$, $v_{term}$ & $S_c$, $C_{D,c}$, $\vec{F}_c$, $V_{term}$ & \texttt{p.recovery.S}, \texttt{p.recovery.Cd}, \texttt{forces.}\allowbreak\texttt{parachute\_}\allowbreak\texttt{drag\_}\allowbreak\texttt{inertial} & canopy area, drag coefficient, drag force, terminal speed & $\mathrm{m^2}$, ---, N, m/s \\
$F_b$ & $A_{brake}$ & \texttt{control.airbrake\_*} & air-brake axial force & N \\
\hline
\multicolumn{5}{l}{\textbf{Analysis}} \\
\hline
$\hat{\omega}=\omega_r c/(2v_{rel})$ & $\hat{q}$ & --- & non-dimensional pitch rate & --- \\
$M_{\theta_t}=(l_g-l_t)\,T\sin\theta_t$ & $M_\delta$ & --- & misalignment disturbance torque & N\,m \\
$s_i$ & $l_i$ & --- & position of a body point relative to the CoG ($C_{m,q}$ derivation) & m \\
$M_\alpha$, $M_q$, $\omega_n$, $\zeta$ & same & --- & short-period parameters & $\mathrm{1/s^2}$, 1/s, rad/s, --- \\
\end{longtable}
}
$A$ is only the nose tip, the reference area being $S$ ($A_e$, $A_f$, $A_m$ carry their own
subscripts). Some letters carry two meanings and are told apart by context only: $g$ is both the
centre of gravity and the gravitational acceleration, and as a subscript both the point ($l_g$,
$r_g$) and the gravity frame ($\mathcal{F}_g$, $\hat{z}_g$); $t$ as a point is likewise distinct
from time $t$, and $h$ is the RK4 step in \autoref{sec:states} but the altitude elsewhere. The
mass model of \autoref{sec:mass} uses the subscript $m$ for the motor and $r$ for the dry rocket,
so $l_p$ is only the aerodynamic reference point; in the code \texttt{p.weight.l\_p} is the
propellant centroid (the motor's $l_{g,m}$).
